\section{Introduction}\label{sec-intro}

\subsection{Contexte}
Une fonction de contrôle indépendant des prix (IPV) contrôle les paramètres que le front office (FO) utilise pour
valoriser et gérer le risque des dérivés actions. Pour le repo, la référence n'est pas observée directement. Elle est
\emph{implicitée} : Totem publie des prix de consensus de calls et de puts, dont on extrait les forwards par la
parité call--put, et un service Totem distinct publie des dividendes de consensus. La courbe de repo qui réconcilie
les deux est ensuite comparée aux repos marqués par le FO, et les écarts sont multipliés par les sensibilités FO pour
produire des lignes de contrôle (écarts de valorisation, provisions, ajustements de valorisation prudente).

Deux caractéristiques de Totem rendent l'exercice délicat. Le service options et le service dividendes sont alimentés
par des contributeurs différents, et dans chaque service le nombre de contributeurs varie avec la maturité ; une action
typique peut compter quinze contributeurs à un an sur les options et dix sur les dividendes, moins au-delà de deux ans.
Le repo implicite mélange les deux moyennes de consensus : il hérite donc du bruit des deux et des artefacts de leur
décalage. Les contrôleurs le constatent sous forme de courbes de repo implicite irrégulières et de sauts injustifiés.

\subsection{Questions traitées}
Cette note répond à quatre questions :
\begin{enumerate}[label=(\alph*)]
\item \textbf{D'où viennent les courbes de repo implicite irrégulières ?} La \cref{sec-anatomy} décompose exactement
le repo implicite en repo du panel options plus un écart de panel, et établit la signature de chaque mécanisme
(écart de panel, sortie de contributeurs, date de détachement, profil de détachement périmé, erreur sur une période de
dividendes, décalage du spot, définition brute/nette).
\item \textbf{Une meilleure méthode statistique peut-elle retrouver le « vrai » repo ?} Non : le \cref{thm-nonid}
montre que les données de consensus ne permettent pas d'identifier la répartition entre repo et dividendes du panel
options.
\item \textbf{Faut-il une courbe de repo Totem propre pour contrôler le repo ?} Pour une exposition linéaire en
forward, non : les \cref{thm-control,thm-exact} démontrent que la somme des lignes dividende et repo ne dépend pas
du panel dividendes, et la \cref{prop-materiality} mesure, à partir des seules sensibilités, l'importance de la
répartition pour un portefeuille donné.
\item \textbf{Comment détecter les courbes aberrantes ?} Les \cref{sec-noise,sec-detection} construisent un
détecteur structurel : un modèle de bruit calculé à partir des propres statistiques de Totem, un gabarit de filtre
adapté par mécanisme, un taux de fausses alertes calibré et une explication en unités naturelles pour chaque alerte.
La \cref{sec-benchmark} le teste, et la \cref{sec-alternatives} explique la place des coefficients B-spline, des
mélanges gaussiens et des petits réseaux de neurones.
\end{enumerate}

\subsection{Principaux résultats}
\begin{insight}
\begin{itemize}[leftmargin=12pt]
\item Au premier ordre, $b_{\text{implicite}}=\bA+(\DA-\DB)/\bigl(T(S-D)\bigr)$ (\cref{prop-gap}). L'écart seul
produit \vGapOneY\,pb de bruit de repo à un an et \vGapFiveY\,pb à cinq ans pour une action typique.
\item Les mécanismes ont des signatures distinctes dans l'\emph{espace de portage} $B=b\,T$ : une constante (spot),
une marche (montant de dividende, changement de panel), un pic (forward aberrant, date de détachement), un motif
saisonnier (profil d'allocation).
\item Pour un portefeuille linéaire en forward, le P\&L Totem en réévaluation complète est \emph{exactement}
indépendant du panel dividendes lorsque les échéances sont des piliers Totem (\cref{thm-exact}) ; la décomposition en
lignes le reproduit au premier ordre (\cref{thm-control}). Le panel dividendes n'intervient qu'à travers
l'exposition non-forward $g=s_D+s_bJ$.
\item Les statistiques de test du détecteur sont des filtres adaptés exacts (\cref{thm-glrt}). Une seule échelle de
variance par catégorie, calibrée sur un historique sain, fixe exactement le taux de fausses alertes (\cref{prop-kappa}).
\item Sur le banc d'essai (\vNTest{} sous-jacents synthétiques de test, chaque détecteur à \vFPR{} de fausses
alertes), le détecteur structurel et un mélange gaussien sur les coefficients B-spline de l'écart Totem--FO détectent
en moyenne une part comparable des anomalies (\vMeanStruct{} et \vMeanGmmDiff). Le détecteur structurel trouve en
outre les forwards aberrants (\vFwdStruct{} contre au plus \vFwdShape), explique correctement \vAttr{} de ses alertes
et ne confond pas une véritable erreur de marquage FO avec une courbe Totem aberrante (\vFoStruct{} contre
\vFoGmmDiff).
\end{itemize}
\end{insight}

\subsection{Guide de lecture}
La \cref{sec-model} fixe les notations et démontre les résultats élémentaires sur lesquels repose la suite
(bootstrap, parité, allocation). La \cref{sec-anatomy} pose le diagnostic ; la \cref{sec-control} établit l'identité
de contrôle ; les \cref{sec-noise,sec-detection} construisent le détecteur ; la \cref{sec-benchmark} apporte les
éléments empiriques ; la \cref{sec-alternatives} compare avec d'autres approches ; la \cref{sec-ops} décrit la mise en
œuvre. Les démonstrations suivent chaque énoncé ; les calculs plus longs figurent en \cref{app-proofs}. Chaque exemple
chiffré est produit par le code du dépôt (\cref{app-repro}).
